% Generated semantic source for AI Almanach v3.
% Edit v3 directly after initial generation; do not overwrite
% editorial changes by rerunning the converter casually.

% ==== Artificial Intelligence in Supply Chain Management =============================================================================
\chapter{Artificial Intelligence in Supply Chain Management}

\chapterlead{How can uncertain demand and constrained resources be coordinated?}{Observe}{Forecast}{Optimise}{Recover}{supply-chain}

% ==== Section 1: Supply Chain Management Fundamentals ===============================================================
\label{ch:supply-chain}

\begin{learningobjectives}
After working through this chapter you should be able to:
\begin{itemize}
  \item compute safety stock from demand variability and a service-level
        target;
  \item explain why forecast accuracy and inventory performance are different
        objectives;
  \item choose a forecast error metric appropriate to intermittent demand;
  \item describe the bullwhip effect and the information that dampens it;
  \item state why a chronological split is mandatory for demand forecasting.
\end{itemize}
\end{learningobjectives}

\begin{plainlanguage}
A supply chain is a long chain of guesses about the future, each one made
before the previous one's outcome is known. Forecasting is the visible part;
the decisions the forecast feeds --- how much to hold, when to reorder, how much
capacity to book --- are where the money is. A forecast that is more accurate
but arrives too late, or in the wrong unit, or without an uncertainty range,
changes nothing.
\end{plainlanguage}

\begin{engineernote}
Do not optimise the forecast. Optimise the decision. A $5\,\%$ improvement in
mean absolute error is worth nothing if the reorder policy rounds to pallet
quantities anyway. Conversely, an unchanged point forecast paired with a
credible \emph{prediction interval} can cut safety stock substantially, because
safety stock is driven by uncertainty rather than by the central estimate.
\end{engineernote}

\section{Supply Chain Management Fundamentals}

% ---- 14.1.1 Supply Chain Concepts -----------------------------------------------------------------------
\subsection{Supply Chain Concepts}


\begin{v3topic}{Supply Chain vs.\ Supply Network}
    A \textbf{supply chain} is the sequence of stages from raw material through production to end customer\textsuperscript{\cite[][pp.~1--4]{chopraSupplyChainManagement2019}}.
    In practice, a firm is embedded in a \textbf{supply network}: multiple upstream tiers (Tier-1, Tier-2, \ldots suppliers) and multiple downstream channels (distributors, retailers, end consumers).
    Supply Chain Management (SCM) coordinates flows of \emph{material}, \emph{information}, and \emph{finance} across this network to maximise total value created.

\end{v3topic}
\begin{v3topic}{SCOR Reference Model}
    The \textbf{Supply Chain Operations Reference (SCOR)} model defines five cross-functional process domains\textsuperscript{\cite[][pp.~5--6]{chopraSupplyChainManagement2019}}:
    \begin{itemize}
        \item \textbf{Plan} --- aggregate demand/supply balancing, network design
        \item \textbf{Source} --- supplier selection, procurement, inbound logistics
        \item \textbf{Make} --- production scheduling, quality, maintenance
        \item \textbf{Deliver} --- order management, warehousing, outbound logistics
        \item \textbf{Return} --- reverse logistics, warranty, recycling
    \end{itemize}

\end{v3topic}
\begin{v3topic}{The Bullwhip Effect}
    The \textbf{bullwhip effect} (Lee et al., 1997) describes how small fluctuations in end-consumer demand amplify progressively as they propagate upstream\textsuperscript{\cite[][pp.~490--493]{chopraSupplyChainManagement2019}}.
    A retailer orders slightly more to be safe; the wholesaler amplifies this; the manufacturer sees enormous swings.
    Root causes include: demand forecast updating, order batching, price fluctuations, and rationing/shortage gaming.
\visualseparator
    \factvisual{
        \textbf{Read the amplification upstream.}
        \begin{itemize}
            \item Customer demand varies modestly.
            \item Retailer orders respond to forecasts and safety stock.
            \item Manufacturer orders inherit delays and larger corrections.
        \end{itemize}
        \textbf{AI mitigation:} shared point-of-sale data, collaborative
        forecasting, and automated replenishment reduce information delay and
        avoid each tier independently over-correcting.
    }{
        \begin{tikzpicture}
            \begin{axis}[
                almanach axis,
                height=50mm,
                xmin=1,xmax=8,
                ymin=25,ymax=95,
                xlabel={period},
                ylabel={relative units},
                xtick={1,2,3,4,5,6,7,8},
                legend style={at={(0.5,1.02)},anchor=south},
                legend columns=3]
                \addplot[AlmanachInk!65,thick,mark=*]
                    coordinates {(1,52)(2,55)(3,50)(4,57)(5,53)(6,58)(7,51)(8,55)};
                \addlegendentry{demand}
                \addplot[AlmanachTeal,thick,mark=square*]
                    coordinates {(1,49)(2,61)(3,43)(4,67)(5,45)(6,70)(7,42)(8,65)};
                \addlegendentry{retailer}
                \addplot[AlmanachOchre,very thick,mark=triangle*]
                    coordinates {(1,42)(2,76)(3,32)(4,86)(5,28)(6,91)(7,31)(8,82)};
                \addlegendentry{factory}
            \end{axis}
        \end{tikzpicture}
        \visualcaption{Illustrative bullwhip pattern: order variability grows
          as delayed, locally optimised decisions move upstream.}
          {fig:bullwhip-amplification}
    }

\end{v3topic}

% ---- 14.1.2 Supply Chain 4.0 ---------------------------------------------------------------------------
\subsection{Supply Chain 4.0}


\begin{v3topic}{Supply Chain 4.0: Digital Transformation}
    Supply Chain 4.0 applies the technologies of Industry 4.0 to logistics and operations\textsuperscript{\cite[][pp.~17--22]{chopraSupplyChainManagement2019}}:
    \begin{itemize}
        \item \textbf{IoT and RFID} for real-time asset and inventory tracking
        \item \textbf{Big data} analytics for demand sensing and supplier scoring
        \item \textbf{AI/ML} for forecasting, routing, and anomaly detection
        \item \textbf{Cloud platforms} for cross-enterprise data sharing
        \item \textbf{Autonomous robots} (AGVs, drones) in warehouses
    \end{itemize}

\end{v3topic}
\begin{v3topic}{Blockchain for Traceability}
    A \textbf{blockchain} is an append-only distributed ledger where transactions are grouped in cryptographically linked blocks\textsuperscript{\cite[][pp.~8--9]{chopraSupplyChainManagement2019}}.
    In supply chains, every handoff (harvest, processing, shipment, customs) is recorded immutably.
    Benefits: anti-counterfeiting, provenance verification (food safety, pharmaceuticals), automated payment via smart contracts, and reduced reconciliation overhead between trade partners.

\end{v3topic}
\begin{v3topic}{Digital Twins in Supply Chains}
    A \textbf{digital twin} is a real-time virtual replica of a physical system: a warehouse, a production line, or an entire network\textsuperscript{\cite[][pp.~2403--2405]{bocewiczDigitalTwinManufacturing2021}}.
    Sensor data continuously updates the model; changes can be tested \emph{in silico} before physical implementation.
    Supply chain use cases: capacity planning under disruption scenarios, predictive maintenance of transport assets, optimisation of warehouse slotting, and network redesign after geopolitical events.
\visualseparator
    \textbf{Workflow}: Physical assets $\xrightarrow{\text{IoT sensors}}$ Digital replica $\xrightarrow{\text{simulation/ML}}$ Decision recommendations $\xrightarrow{\text{feedback}}$ Physical assets.

\end{v3topic}

% ---- 14.1.3 AI Techniques for Supply Chains ------------------------------------------------------------
\subsection{AI Techniques for Supply Chains}


\begin{v3topic}{AI Technique Overview}
    \begin{tblr}{width=\linewidth,colspec={Q[l,m,3cm]X X},hlines,vlines}
        \hline
        \textbf{Technique}           & \textbf{Core idea}                                    & \textbf{SC Application}                           \\
        \hline
        ML / Statistical             & Learn patterns from historical data                   & Demand forecasting, quality prediction             \\
        \hline
        Optimisation (LP/MIP)        & Minimise/maximise objective subject to constraints    & Network design, inventory, routing                 \\
        \hline
        Reinforcement Learning       & Policy learned via environment interaction            & Dynamic routing, replenishment, scheduling         \\
        \hline
        RPA                          & Automate rule-based digital workflows                 & Purchase order creation, invoice processing        \\
        \hline
        Multi-Agent Systems          & Distributed agents coordinate autonomously            & Decentralised supplier negotiation, fleet dispatch \\
        \hline
    \end{tblr}
\captionof{table}{Machine-learning techniques mapped to supply-chain applications.}
\label{tab:sc-techniques}

\end{v3topic}

% ==== Section 2: Demand Forecasting ===========================================================================
\section{Demand Forecasting}

% ---- 14.2.1 Traditional Methods -------------------------------------------------------------------------
\subsection{Traditional Methods}


\begin{v3topic}{Exponential Smoothing (SES)}
    \textbf{Simple Exponential Smoothing} weights recent observations more heavily via smoothing parameter $\alpha \in (0,1)$\textsuperscript{\cite[][pp.~75--77]{hyndmanForecastingPrinciples2021}}:
    \begin{equation}
        \hat{y}_{t+1} = \alpha\, y_t + (1-\alpha)\,\hat{y}_t
    \end{equation}
    Equivalently, $\hat{y}_{t+1} = \alpha \sum_{j=0}^{\infty}(1-\alpha)^j y_{t-j}$ --- an exponentially decaying weighted average.
    Suitable for data with no trend or seasonality. A large $\alpha$ makes the forecast react quickly but noisily; a small $\alpha$ produces smooth but slow-to-adapt forecasts.

\end{v3topic}
\begin{v3topic}{Holt-Winters / ETS}
    \textbf{Holt's linear method} adds a trend component $b_t$; \textbf{Holt-Winters} further adds seasonality $s_t$\textsuperscript{\cite[][pp.~82--95]{hyndmanForecastingPrinciples2021}}:
    \begin{equation}
        \hat{y}_{t+h} = \ell_t + h\,b_t + s_{t+h-m}
    \end{equation}
    where $\ell_t$ is level, $b_t$ trend, $s_t$ seasonal index, $m$ seasonal period.
    The \textbf{ETS} (Error--Trend--Seasonality) framework unifies exponential smoothing models through state-space representations, enabling likelihood-based parameter estimation and model selection via AIC.

\end{v3topic}
\begin{v3topic}{ARIMA Models}
    An \textbf{ARIMA}$(p,d,q)$ model combines three components\textsuperscript{\cite[][pp.~233--257]{hyndmanForecastingPrinciples2021}}:
    \begin{itemize}
        \item \textbf{AR}$(p)$: $y_t = \phi_1 y_{t-1}+\cdots+\phi_p y_{t-p} + \varepsilon_t$
        \item \textbf{I}$(d)$: $d$-th order differencing to achieve stationarity
        \item \textbf{MA}$(q)$: $y_t = \varepsilon_t + \theta_1\varepsilon_{t-1}+\cdots+\theta_q\varepsilon_{t-q}$
    \end{itemize}
    Model order is selected with ACF/PACF plots or auto-ARIMA (AIC/BIC minimisation).
    SARIMA$((p,d,q)(P,D,Q)_m)$ extends this to seasonal data.

\end{v3topic}
\begin{v3topic}{State-Space and Bayesian Structural Time Series}
    \textbf{State-space models} represent time series via a latent state vector $\boldsymbol{\alpha}_t$ evolving according to a transition equation, with observations linked via a measurement equation\textsuperscript{\cite[][pp.~293--310]{hyndmanForecastingPrinciples2021}}.
    The Kalman filter provides optimal linear prediction.
    \textbf{Bayesian Structural Time Series (BSTS)} add spike-and-slab priors to select relevant covariates and decompose series into interpretable components (trend, seasonality, regression), quantifying forecast uncertainty through the posterior distribution.

\end{v3topic}

% ---- 14.2.2 ML-Based Forecasting ------------------------------------------------------------------------
\subsection{ML-Based Forecasting}


\begin{v3topic}{LSTMs for Time Series}
    \textbf{Long Short-Term Memory (LSTM)} networks address the vanishing gradient problem of plain RNNs through gated memory cells\textsuperscript{\cite[][pp.~1748--1751]{limTemporalFusionTransformers2021}}:
    \begin{itemize}
        \item \textbf{Forget gate} $f_t$: discard irrelevant past state
        \item \textbf{Input gate} $i_t$: add new information to cell state $c_t$
        \item \textbf{Output gate} $o_t$: expose cell state as hidden state $h_t$
    \end{itemize}
    LSTMs capture long-range temporal dependencies and are effective for univariate and multivariate demand forecasting, outperforming ARIMA when series exhibit complex nonlinear patterns.

\end{v3topic}
\begin{v3topic}{Gradient Boosting for Forecasting}
    Tree-based gradient boosting methods (\textbf{XGBoost}, \textbf{LightGBM}) have become the workhorse for tabular demand forecasting in retail and supply chain competitions\textsuperscript{\cite[][p.~905]{huberDataDrivenNewsvendor2019}}.
    Key advantages: handle large feature sets (calendar effects, price, promotions, product attributes), fast training, interpretable feature importance, and no assumption of stationarity.
    The model is a sum of regression trees: $\hat{y}_t = \sum_{k=1}^K f_k(\mathbf{x}_t)$, trained with additive stage-wise optimisation.

\end{v3topic}
\begin{v3topic}{Temporal Fusion Transformer (TFT)}
    The \textbf{Temporal Fusion Transformer} combines recurrent layers, multi-head attention, and variable selection networks for interpretable multi-horizon forecasting\textsuperscript{\cite[][pp.~1748--1764]{limTemporalFusionTransformers2021}}:
    \begin{itemize}
        \item \textbf{Variable Selection Network}: gates out irrelevant covariates
        \item \textbf{LSTM encoder-decoder}: encodes observed + future known inputs
        \item \textbf{Temporal self-attention}: identifies salient past time steps
        \item \textbf{Quantile regression}: outputs prediction intervals
    \end{itemize}
    TFT achieves state-of-the-art accuracy on multiple benchmark datasets.

\end{v3topic}
\begin{v3topic}{Probabilistic Forecasting}
    Point forecasts alone are insufficient for inventory decisions --- we need the \textbf{full demand distribution}\textsuperscript{\cite[][p.~37]{wickDemandForecastingIndividual2021}}.
    Probabilistic forecasting methods:
    \begin{itemize}
        \item Quantile regression: model arbitrary quantiles of $P(y_{t+h} \le q)$
        \item Normalising flows: learn flexible joint distributions over multivariate series\textsuperscript{\cite{rasulMultivariateProbabilisticTime2021}}
        \item Conformal prediction: distribution-free coverage guarantees
        \item Bayesian neural networks: posterior predictive intervals
    \end{itemize}
    The output is used directly in inventory optimisation (e.g.\ to compute safety stock or order-up-to levels at a given service level).

\end{v3topic}

% ---- 14.2.3 The Newsvendor Problem -----------------------------------------------------------------------
\subsection{The Newsvendor Problem}


\begin{v3topic}{Classic Single-Period Newsvendor}
    The \textbf{newsvendor problem} is a canonical single-period inventory model under stochastic demand $D$\textsuperscript{\cite[][pp.~537--541]{hillierIntroductionOperationsResearch2020}}.
    Order quantity $Q$ must be committed before demand is realised.
    The cost function (with unit overage cost $c_o = c - s$ and underage cost $c_u = p - c$) is:
    \begin{equation}
        \mathbb{E}[C(Q)] = c_o\,\mathbb{E}[(Q-D)^+] + c_u\,\mathbb{E}[(D-Q)^+]
    \end{equation}
    The optimal order quantity satisfies the \textbf{critical ratio} condition:
    \begin{equation}
        F(Q^*) = \frac{c_u}{c_u + c_o} = \text{CR}
    \end{equation}
    i.e.\ $Q^*$ equals the $\text{CR}$-th quantile of the demand distribution $F$.

\end{v3topic}
\begin{v3topic}{Demand Models and Censoring}
    Demand is modelled as a random variable $D \sim F_\theta$. Common parametric families: Normal $\mathcal{N}(\mu,\sigma^2)$, Poisson for count data, Negative Binomial for overdispersed counts, and log-normal for positive skewed demand\textsuperscript{\cite[][p.~540]{hillierIntroductionOperationsResearch2020}}.
    \textbf{Censored demand}: when stock-outs occur, actual demand exceeds sales data. The EM algorithm or maximum likelihood with censoring adjustments recovers true demand parameters\textsuperscript{\cite[][pp.~490--495]{porteusFundamentationsStochasticInventory2002}}.
    Ignoring censoring \emph{systematically underestimates} mean demand.

\end{v3topic}
\begin{v3topic}{Data-Driven Newsvendor (Ban \& Rudin)}
    \textbf{Ban and Rudin (2019)} replace distributional assumptions with a direct data-driven approach: given feature vector $\mathbf{x}_t$ (price, calendar, promotions) and past demand observations $\{(\mathbf{x}_i, d_i)\}$, choose $Q$ to minimise empirical expected cost using sample-average approximation with regularisation\textsuperscript{\cite[][pp.~90--108]{banBigDataNewsvendor2019}}.
    Key insight: the optimal policy $Q^*(\mathbf{x})$ can be estimated via \textbf{weighted empirical quantile regression} --- nonparametrically consistent under mild conditions.
    \textbf{Bertsimas and Kallus (2020)} extend this ``predict-then-optimise'' to general stochastic programs via \emph{prescriptive analytics}\textsuperscript{\cite[][pp.~1025--1044]{bertsimasPredictivePrescriptiveAnalytics2020}}.

\end{v3topic}

% ==== Section 3: Inventory Management ===========================================================================
\section{Inventory Management}

% ---- 14.3.1 Economic Order Quantity -------------------------------------------------------------------------
\subsection{Economic Order Quantity (EOQ)}


\begin{v3topic}{Harris EOQ Formula}
    The \textbf{Economic Order Quantity} model (Harris, 1913) minimises the sum of ordering and holding costs for a single item with constant demand rate $\lambda$\textsuperscript{\cite{harrisHowManyParts1913}}:
    \begin{equation}
        TC(Q) = \frac{\lambda}{Q}\,K + \frac{Q}{2}\,h
    \end{equation}
    where $K$ = fixed order cost, $h = ic$ = unit holding cost ($i$ = inventory carrying rate, $c$ = unit purchase cost).
    Setting $\frac{d\,TC}{dQ}=0$ yields:
    \begin{equation}
        Q^* = \sqrt{\frac{2\lambda K}{h}}
    \end{equation}
    The EOQ is remarkably robust: a 30\% error in $K$ or $h$ leads to only a $\sim$14\% error in $Q^*$ (``square root rule'').

\end{v3topic}
\begin{v3topic}{EOQ Assumptions and Extensions}
    The classic EOQ assumes\textsuperscript{\cite[][pp.~936--942]{hillierIntroductionOperationsResearch2020}}:
    \begin{itemize}
        \item Constant, deterministic demand rate $\lambda$
        \item Instantaneous replenishment (no lead time)
        \item No shortages allowed
        \item Single item, no quantity discounts
    \end{itemize}
    \textbf{Extensions}: EOQ with quantity discounts, EOQ with backordering, production order quantity (POQ) for finite production rate, joint replenishment for multiple items with shared ordering costs, and stochastic demand extensions leading to $(s,Q)$ policies.

\end{v3topic}

% ---- 14.3.2 Inventory Policies -------------------------------------------------------------------------------
\subsection{Inventory Policies and Service Levels}


\begin{v3topic}{Continuous Review $(s{,}Q)$ Policy}
    Under \textbf{continuous review}, inventory position is monitored continuously\textsuperscript{\cite[][pp.~362--384]{galliherDynamicsContinuousReview1959}}.
    The $(s,Q)$ (reorder point, order quantity) policy: when inventory position drops to reorder point $s$, order a fixed quantity $Q$.
    The reorder point protects against demand during lead time $L$:
    \begin{equation}
        s = \mu_L + z_\alpha \sigma_L
    \end{equation}
    where $\mu_L = \lambda L$ is mean lead-time demand, $\sigma_L = \sigma\sqrt{L}$ its standard deviation, and $z_\alpha$ is the standard normal quantile corresponding to desired service level.

\end{v3topic}
\begin{v3topic}{Periodic Review $(s{,}S)$ Policy}
    In the \textbf{periodic review} setting, inventory is checked every $T$ periods\textsuperscript{\cite[][p.~2]{porteusFundamentationsStochasticInventory2002}}.
    The $(s,S)$ policy: if inventory falls below $s$, order up to $S$ (``order-up-to'' level).
    Optimal under fixed ordering costs (Scarf, 1959\textsuperscript{\cite{scarfBayesSolutionsStatistical1959}}).
    The gap $S - s$ balances setup cost $K$ against holding and shortage costs.

\end{v3topic}
\begin{v3topic}{Service Levels and Safety Stock}
    Two common service level
    definitions\textsuperscript{\cite[][pp.~300--305]{vandeputInventoryOptimization2020}}:
    \par\smallskip\noindent
    \begin{tblr}{width=\linewidth,colspec={Q[l,m,3cm]X X},hlines,vlines}
        \textbf{Metric}       & \textbf{Definition}                                        & \textbf{Typical target} \\
        Cycle service level (\textbf{Type~I})   & $P(\text{no stockout per cycle}) = F(s)$      & 90--99\% \\
        Fill rate (\textbf{Type~II})             & Fraction of demand satisfied from stock       & 95--99.5\% \\
    \end{tblr}
\captionof{table}{Service-level definitions and typical targets. The relationship to safety stock is convex, so each additional nine costs disproportionately more inventory.}
\label{tab:sc-service-levels}
    \textbf{Safety stock} (SS) is the buffer held against demand and lead-time variability:
    $\text{SS} = z_\alpha \sqrt{L \sigma_d^2 + \mu_d^2 \sigma_L^2}$
    where $\sigma_d$ = demand std.\ dev., $\sigma_L$ = lead-time std.\ dev.

\end{v3topic}

% ---- 14.3.3 Advanced Inventory Topics -----------------------------------------------------------------------
\subsection{Advanced Inventory Topics}


\begin{v3topic}{Perishable Inventory}
    For products that expire after a fixed lifetime (food, pharmaceuticals, blood), classical models must be extended\textsuperscript{\cite{nahmiasOptimalOrderingPolicies1973}}.
    Key challenges: first-in-first-out (FIFO) or last-in-first-out (LIFO) issuing policies affect freshness and waste; outdating costs must be balanced against shortage costs.
    ML approaches train neural policies that incorporate remaining shelf-life as state variable, outperforming static order-up-to policies in simulation studies.

\end{v3topic}
\begin{v3topic}{Minimum Order Quantities and Multi-Echelon}
    \textbf{Minimum order quantities (MOQ)} imposed by suppliers create discrete, non-convex feasible sets, invalidating the standard EOQ formula.
    Heuristic approaches: round-up to the nearest MOQ multiple, simulate total cost sensitivity, or solve as an integer programming problem\textsuperscript{\cite[][pp.~210--215]{vandeputInventoryOptimization2020}}.
    \textbf{Multi-echelon inventory} coordinates stock across distribution centres, regional warehouses, and retail outlets; Clark and Scarf (1960) showed that echelon stock formulations decompose the problem across echelons for serial systems, reducing computational complexity.

\end{v3topic}

% ==== Section 4: Strategic and Operational AI Applications ==============================================
\section{Strategic and Operational AI Applications}

% ---- 14.4.1 Supply Chain Network Design -------------------------------------------------------------------
\subsection{Supply Chain Network Design}


\begin{v3topic}{Facility Location Problem}
    Network design determines the number, location, and capacity of facilities (plants, DCs, warehouses) to minimise total cost (fixed + variable operating + transportation)\textsuperscript{\cite[][pp.~146--160]{chopraSupplyChainManagement2019}}.
    The \textbf{uncapacitated facility location} integer program:
    \begin{equation}
        \min \sum_j f_j y_j + \sum_{i,j} c_{ij} x_{ij}
        \quad \text{s.t.}\ \sum_j x_{ij}=1,\ x_{ij}\le y_j,\ y_j\in\{0,1\}
    \end{equation}
    where $y_j$ = open/close binary, $x_{ij}$ = fraction of customer $i$ served by facility $j$.
    Solved with branch-and-bound (MIP solvers) or Benders decomposition for large instances.

\end{v3topic}
\begin{v3topic}{Scenario Planning and Resilience}
    Deterministic network designs can fail under disruptions (natural disasters, geopolitical risk, pandemics)\textsuperscript{\cite[][pp.~600--605]{chopraSupplyChainManagement2019}}.
    \textbf{Robust optimisation} designs networks that perform well under worst-case demand/disruption scenarios; \textbf{stochastic programming} incorporates scenario trees with probabilities.
    ML enhances resilience by: predicting supplier failure probabilities from financial and ESG signals, flagging concentration risk (geographic or single-source), and recommending diversification.

\end{v3topic}

% ---- 14.4.2 Supplier Selection and Procurement -----------------------------------------------------------
\subsection{Supplier Selection and Procurement}


\begin{v3topic}{Multi-Criteria Supplier Scoring}
    Supplier selection integrates quantitative (price, quality, delivery reliability, lead time) and qualitative (sustainability, financial stability, geographic risk) criteria\textsuperscript{\cite[][pp.~308--314]{chopraSupplyChainManagement2019}}.
    Classical methods: Analytic Hierarchy Process (AHP) for pairwise criteria weighting, TOPSIS for ranking alternatives.
    ML approaches: gradient boosted trees trained on historical supplier performance and quality control data, producing a real-time supplier risk score.
    Feature engineering: on-time delivery rate, defect rate, capacity utilisation, ESG ratings, financial ratios.

\end{v3topic}
\begin{v3topic}{Spend Analytics}
    \textbf{Spend analytics} aggregates and classifies procurement transactions to identify savings opportunities and consolidate the supplier base\textsuperscript{\cite[][pp.~316--320]{chopraSupplyChainManagement2019}}.
    Key ML tasks: natural language classification of invoice descriptions into UNSPSC/eCl@ss taxonomy (typically fine-tuned BERT or logistic regression on TF-IDF features); duplicate vendor detection (entity resolution); price benchmarking against market indices.
    Pareto analysis (80\% of spend in 20\% of categories) directs strategic sourcing effort.

\end{v3topic}

% ---- 14.4.3 Route Optimisation -------------------------------------------------------------------------------
\subsection{Route Optimisation}


\begin{v3topic}{Vehicle Routing Problem (VRP)}
    The \textbf{Capacitated VRP} seeks minimum-cost routes for a fleet of vehicles, each with capacity $C$, to serve customers with demands $d_i$ from a depot\textsuperscript{\cite[][pp.~420--440]{hillierIntroductionOperationsResearch2020}}.
    It generalises the \textbf{Travelling Salesman Problem (TSP)}, which is NP-hard.
    Solution approaches: Clarke-Wright savings heuristic, LKH meta-heuristic for TSP, branch-price-and-cut exact MIP solvers, and \textbf{attention-based neural combinatorial optimisation} (pointer networks, transformer-based).

\end{v3topic}
\begin{v3topic}{Dynamic Routing with Reinforcement Learning}
    Real-world routing involves stochastic, time-varying conditions (traffic, dynamic orders, vehicle breakdowns)\textsuperscript{\cite[][pp.~1--3]{panJobShopSchedulingRL2021}}.
    \textbf{Deep RL (DRL)} learns routing policies directly from environment interaction, representing state (vehicle positions, remaining deliveries, time windows) and actions (next customer to visit).
    Actor-Critic architectures (PPO, A3C) outperform static VRP heuristics in simulations of dynamic last-mile delivery, adapting in real time without re-solving an NP-hard combinatorial problem.

\end{v3topic}
\begin{v3topic}{Last-Mile Delivery}
    Last-mile delivery accounts for 40--60\% of total logistics cost\textsuperscript{\cite[][pp.~444--448]{chopraSupplyChainManagement2019}}.
    AI applications:
    \begin{itemize}
        \item \textbf{Delivery time prediction}: regression models (gradient boosting) on driver, route, and parcel features for ETA estimation
        \item \textbf{Stop sequencing}: real-time reoptimisation with learned heuristics
        \item \textbf{Drone/autonomous vehicle routing}: energy-constrained routing with RL
        \item \textbf{Crowdsourced delivery}: matching demand (parcels) to supply (couriers) via two-sided marketplace algorithms
        \item \textbf{Micro-fulfillment centres}: ML-driven slotting to minimise pick travel distance
    \end{itemize}

\end{v3topic}

% ---- 14.4.4 Sales and Operations Planning ---------------------------------------------------------------
\subsection{Sales and Operations Planning (S\&OP)}


\begin{v3topic}{The S\&OP Process}
    \textbf{Sales and Operations Planning} is the monthly cross-functional process that aligns demand forecasts with supply capacity to achieve business targets\textsuperscript{\cite[][pp.~188--196]{chopraSupplyChainManagement2019}}.
    The five-step process:
    \begin{enumerate}
        \item Product review (new/discontinued items)
        \item Demand review (statistical baseline + commercial adjustments)
        \item Supply review (capacity and constraints)
        \item Pre-S\&OP (gap analysis, scenario options)
        \item Executive S\&OP (decision and sign-off)
    \end{enumerate}

\end{v3topic}
\begin{v3topic}{AI in S\&OP}
    ML transforms the S\&OP process from a backward-looking spreadsheet exercise to a forward-looking decision-support system\textsuperscript{\cite[][pp.~196--201]{chopraSupplyChainManagement2019}}:
    \begin{itemize}
        \item \textbf{Consensus forecasting}: hierarchical time-series models reconcile forecasts across product/geography hierarchies (top-down, bottom-up, middle-out)
        \item \textbf{What-if simulation}: digital twin evaluates supply plans under demand scenarios
        \item \textbf{Automated alerts}: ML flags gaps between demand plan and constrained supply before the S\&OP meeting
        \item \textbf{Bias detection}: statistical tests identify systematic forecast over- or under-bias by planner, region, or product family
    \end{itemize}

\end{v3topic}

% ==== Section 5: Transparency and Risk ======================================================================
\section{Transparency and Risk}

% ---- 14.5.1 Supply Chain Visibility -----------------------------------------------------------------------
\subsection{Supply Chain Visibility}


\begin{v3topic}{Real-Time Tracking and Order Prediction}
    \textbf{End-to-end visibility} requires integrating data from ERP systems, carrier APIs, IoT sensors, and partner portals into a unified tracking platform\textsuperscript{\cite[][pp.~72--80]{chopraSupplyChainManagement2019}}.
    ML applications:
    \begin{itemize}
        \item \textbf{Predictive ETA}: gradient boosting on shipment history, weather, carrier performance
        \item \textbf{Order peak prediction}: LSTM or Prophet models for order volume spikes (promotional events, seasonality)
        \item \textbf{Exception management}: anomaly detection to flag shipments at risk of delay before they miss SLA
    \end{itemize}

\end{v3topic}
\begin{v3topic}{Customer and Churn Analytics}
    \textbf{Churn prediction} identifies B2B customers at risk of reducing order volume or switching suppliers\textsuperscript{\cite[][p.~85]{chopraSupplyChainManagement2019}}.
    Features: order frequency decline, invoice dispute rate, NPS scores, service failures, competitor pricing.
    Model: gradient boosted classifier or survival model (time-to-churn); SHAP values explain predictions to account managers.
    \textbf{Customer segmentation} via RFM (Recency, Frequency, Monetary value) analysis guides differentiated service levels and inventory positioning.

\end{v3topic}

% ---- 14.5.2 Fraud and Risk Detection -----------------------------------------------------------------------
\subsection{Fraud and Risk Detection}


\begin{v3topic}{Anomaly Detection in Procurement}
    Procurement fraud includes fictitious vendors, split invoices (under approval thresholds), and price collusion.
    ML approaches\textsuperscript{\cite[][pp.~8--10]{chopraSupplyChainManagement2019}}:
    \begin{itemize}
        \item \textbf{Isolation Forest} and \textbf{Autoencoder}: unsupervised detection of unusual transaction patterns
        \item \textbf{Network analysis}: graph-based detection of vendor-employee relationships (conflicts of interest)
        \item \textbf{Benford's Law}: distribution of first digits in legitimate financial data follows a specific law; deviations signal manipulation
        \item \textbf{NLP}: detect suspicious text in contracts, vendor names, and purchase descriptions
    \end{itemize}

\end{v3topic}
\begin{v3topic}{Counterfeit Detection and Supply Chain Security}
    Counterfeit components are a critical risk in electronics, pharmaceuticals, and aerospace supply chains.
    Computer vision (CNN-based classifiers) inspects product markings, holograms, and packaging for deviations from authentic reference images\textsuperscript{\cite{laviollaCNNDefectDetection2019}}.
    Blockchain-anchored digital product passports link physical items (via QR/RFID) to immutable provenance records, making counterfeiting economically infeasible.
    Serialisation and track-and-trace mandates (EU Falsified Medicines Directive, US DSCSA) drive adoption.

\end{v3topic}

% ---- 14.5.3 Challenges --------------------------------------------------------------------------------------
\subsection{Challenges of Applying AI in Supply Chains}


\begin{v3topic}{Data Quality and Organisational Challenges}
    The primary barrier to AI adoption in supply chains is data quality rather than algorithmic sophistication\textsuperscript{\cite[][pp.~620--625]{chopraSupplyChainManagement2019}}:
    \begin{itemize}
        \item \textbf{Data silos}: ERP, WMS, TMS, supplier portals use incompatible schemas; data integration requires significant engineering effort
        \item \textbf{Missing and censored data}: stockouts censor demand; unrecorded returns bias historical data
        \item \textbf{Organisational inertia}: planners distrust black-box AI recommendations; change management and explainability are critical
        \item \textbf{Short history}: new products have little sales history; cold-start forecasting relies on attribute-based (similar item) approaches
    \end{itemize}

\end{v3topic}
\begin{v3topic}{Trust, Accountability, and Accessibility}
    AI deployment in supply chains raises governance questions\textsuperscript{\cite[][pp.~625--630]{chopraSupplyChainManagement2019}}:
    \begin{itemize}
        \item \textbf{Trust}: planners need to understand and override AI recommendations; human-in-the-loop designs maintain planner authority while reducing manual work
        \item \textbf{Accountability}: when an AI-driven procurement decision causes a shortage, responsibility is diffuse; audit trails and decision logging are essential
        \item \textbf{Accessibility}: small and medium enterprises lack the data science talent and IT infrastructure for advanced ML; cloud-based SaaS SCM tools (e.g.\ o9, Kinaxis, Blue Yonder) are democratising access
        \item \textbf{Algorithmic bias}: models trained on pre-COVID data failed during the pandemic; continual monitoring and retraining protocols are mandatory
    \end{itemize}

\end{v3topic}

% ==== Section 6: Further Reading ===========================================================================
%
\section{Deep Dive: Forecast Distributions to Robust Decisions}
\begin{v3topic}{Evaluate at the Decision Horizon}
Use rolling-origin evaluation with the forecast horizons, hierarchy, and data
latency of the planning process.
Report scale-appropriate errors and interval coverage; one aggregate point
metric can hide poor service for intermittent or high-value items.

\end{v3topic}
\begin{v3topic}{Reconcile the Hierarchy}
Product, region, channel, and total forecasts should add consistently.
Reconciliation trades local fit for coherent planning and must respect which
levels contain reliable signal and which are management aggregates.

\end{v3topic}
\begin{v3topic}{Optimise Across Scenarios}
Pass a forecast distribution or scenarios into inventory, routing, or capacity
decisions rather than treating a point forecast as certain demand.
Stress lead time, supplier failure, correlated disruption, and policy delay.
Compare the learned policy with a transparent robust or stochastic-optimisation
baseline.
\end{v3topic}


\begin{workedexample}[Safety stock, and what a service level costs]
Weekly demand averages $\mu = 100$ units with standard deviation
$\sigma = 20$. Lead time is $L = 4$ weeks. Target service level $95\,\%$
($z = 1.65$).

\emph{Demand over the lead time.} Assuming independent weeks, variability adds
in quadrature, not linearly:
\[
  \sigma_L = \sigma\sqrt{L} = 20\sqrt{4} = 40 \text{ units.}
\]

\emph{Safety stock and reorder point.}
\[
  SS = z\,\sigma_L = 1.65\times40 = 66 \text{ units},
  \qquad
  ROP = \mu L + SS = 400 + 66 = 466 \text{ units.}
\]

\emph{Now raise the service level to $99\,\%$} ($z = 2.33$):
\[
  SS = 2.33\times40 = 93 \text{ units.}
\]

\emph{Read the trade.} Four extra percentage points of service cost $27$ extra
units --- a $41\,\%$ increase in safety stock. The relationship is convex: each
additional nine of service costs disproportionately more inventory, and
$100\,\%$ is unreachable at any finite stock. The $\sqrt{L}$ term also explains
why halving lead time is usually worth more than improving the forecast: it
reduces $\sigma_L$ by $29\,\%$ directly.
\end{workedexample}

\begin{cautionbox}[MAPE fails exactly where demand is hardest]
Mean absolute percentage error divides by the actual value, so a week with zero
demand makes it undefined and a week with demand of $1$ makes it explode.
Intermittent demand --- spare parts, slow movers, long-tail SKUs --- is precisely
where this happens, and it is precisely where forecasting is hardest. MAPE also
penalises over-forecasting more heavily than under-forecasting, which quietly
biases a model toward stockouts.

Use MASE, which scales by a naive seasonal baseline and stays defined at zero,
or evaluate the decision directly in units and euros. And always report the
naive baseline --- ``last week's value'' beats a surprising number of deployed
forecasting systems.
\end{cautionbox}

\begin{practicallab}[Forecast, then cost the forecast]
\textbf{Goal.} Show that forecast accuracy and business value are different
things.

\textbf{Protocol.}
\begin{enumerate}
  \item Take any public demand series. Split \emph{chronologically}.
  \item Establish two naive baselines: last value, and same period last year.
  \item Fit one statistical and one machine-learning forecaster. Report MASE
        for all four.
  \item Now convert each forecast into a reorder decision using the safety
        stock formula above, at a $95\,\%$ target.
  \item Simulate the resulting inventory: units held, stockouts, and total cost
        with holding and shortage costs you choose.
  \item Rank the four methods by MASE and by cost. Do the rankings agree?
\end{enumerate}

\textbf{Deliverable.} Two ranked tables and one sentence on any disagreement
between them. \textbf{The point:} step 6 disagrees more often than people
expect, and when it does, the cost ranking is the one that matters.
\end{practicallab}

\begin{checkpoint}
Your forecast MASE improves from $0.95$ to $0.88$. Give two reasons this might
produce no inventory saving at all, and name the change that would.
\end{checkpoint}

\chapterreview{Observe}{Forecast}{Optimise}{Recover}

\section{Further Reading}

\begin{v3topic}{Recommended Sources for AI in Supply Chain Management}
    \begin{itemize}
        \item \fullcite{chopraSupplyChainManagement2019} --- The standard reference for supply chain strategy, planning, and operations. Covers network design, demand management, inventory, transportation, and coordination across all planning horizons.
        \item \fullcite{vandeputInventoryOptimization2020} --- Practical, Python-based guide to inventory models and simulation: EOQ extensions, safety stock, service levels, and data-driven approaches. Essential for practitioners.
        \item \fullcite{hyndmanForecastingPrinciples2021} --- Comprehensive, freely available textbook on statistical and ML forecasting: ETS, ARIMA, dynamic regression, hierarchical forecasting, and neural network methods with R examples.
        \item \fullcite{hillierIntroductionOperationsResearch2020} --- Encyclopedic OR textbook covering linear programming, integer programming, network flows, inventory models, and queuing theory --- the mathematical backbone of supply chain optimisation.
        \item \fullcite{banBigDataNewsvendor2019} --- Seminal paper showing how machine learning can directly optimise newsvendor decisions from data, bypassing the estimate-then-optimise paradigm.
        \item \fullcite{bertsimasPredictivePrescriptiveAnalytics2020} --- Extends the data-driven optimisation framework to general stochastic programs: a unifying treatment of predictive and prescriptive analytics.
        \item \fullcite{limTemporalFusionTransformers2021} --- Introduces the Temporal Fusion Transformer for interpretable multi-horizon forecasting; state-of-the-art on multiple benchmark datasets with attention-based explanations.
        \item \fullcite{peixeiroTimeSeriesForecasting2026} --- Modern guide to foundation models for time series: TimeGPT, Chronos, and fine-tuning strategies for demand forecasting without large training datasets.
    \end{itemize}
\end{v3topic}
